% higgs_composite_qg.tex -- Regular Article for Physical Review D (REVTeX 4.2), revised hypotheses.
% PROJECT spacetimes, folder Revision_SM_Higgs/paper_prd.
% This file is generated from src/higgs_composite_qg.template.tex by scripts/build_paper.py, which inserts every number
% that a program computes, from the data files and the reports of the checks of the folder data/ (the bounds of the
% agreement of two computations rounded up).  Typed in the template are only the parameters of the model and of the
% numerical method, mathematical constants, the quantum numbers that the hypotheses assign, and measured values of the
% Standard Model with their references.
% Build:  python3 paper_prd/scripts/build_paper.py   (writes higgs_composite_qg.tex and runs pdflatex and bibtex)
\documentclass[aps,prd,twocolumn,superscriptaddress,nofootinbib,floatfix,longbibliography]{revtex4-2}
\usepackage{amsmath,amssymb}
\usepackage{graphicx}
\usepackage{bm}
\usepackage[colorlinks=true,linkcolor=blue,citecolor=blue,urlcolor=blue]{hyperref}
\hypersetup{pdftitle={Higgs as a Composite, and Quantum Gravity},pdfauthor={Patrick Lee Nash}}

% the suffixes of the symbols of PROJECT spacetimes: LL (the Lorentz super-loop, Ref. [Nash2026]), SF (the superfluid of
% the earlier hypotheses) and SH (the objects of the revised hypotheses)
\newcommand{\LL}{{\rm LL}}
\newcommand{\SF}{{\rm SF}}
\newcommand{\SH}{{\rm SH}}
\newcommand{\tr}{\mathrm{tr}}

\begin{document}

\title{Higgs as a Composite, and Quantum Gravity}

\author{Patrick Lee Nash, Ph.D.}
\affiliation{PROJECT spacetimes}

\date{November 22, 2026}

\begin{abstract}
The unit bi-octonions form a Moufang loop, the Lorentz super-loop, that contains $SL(2,\mathbb C)$ and eight spinor
charges, carried by an isodoublet $\Psi$ of Dirac spinors. Two members of one isodoublet pair into isospin triplets, so a
superfluid of $\Psi$ alone cannot be the Higgs field of the Standard Model. We give $\Psi$ the hypercharge $\frac14$, add two
Dirac singlets $\chi$ and $\chi'$ with the hypercharges $\frac14$ and $-\frac34$, and let a contact interaction bind a
member of $\Psi$ to a singlet in the spin-zero $s$-wave channel. The two pair fields form a bidoublet of
$SU(2)_L\times SU(2)_R$ with the quantum numbers of two Higgs doublets. In five mean-field methods
(Hubbard-Stratonovich with three decouplings, Kohn-Sham, and Cornwall-Jackiw-Tomboulis) the ground state is the
custodial vacuum: it breaks $SU(2)_L\times U(1)_Y$ to the $U(1)$ of the electric charge, keeps the photon massless and
gives $\rho=1$. Setting its scale to $v=246.22$~GeV puts the momentum cutoff at {{LAM2_TEV}}~TeV. The amplitude mode,
the candidate for the Higgs boson, is a composite scalar bound below the threshold of two quasi-particles. In the
methods in which only the pairing acts, the linear response has the closed form of the theory of Bardeen, Cooper and
Schrieffer and places the mode exactly at the pair self-energy: $M_H/v={{MH_RPA_V}}$, {{FACTOR_RPA}} times the measured
$m_H/v={{MH_SM_V}}$, and over a scan of the coupling from 1.02 to 3 times its critical value, with two masses of the
fields, the ratio is at least {{SCAN_MIN2}}. A
custodial triplet of massive scalars and a Goldstone boson of the particle number come with it. The time-dependent
mean field confirms these frequencies; in the methods with exchange self-energies the mode lies above the threshold and
decays. The superfluid therefore has the quantum numbers of the Higgs field and gives the $W$ and $Z$ bosons their
masses, but in the mean field its amplitude mode is far too heavy. We outline how gravitation will enter through
E.~A.~Lord's gauge theory of the local change of spinor representation; it is not computed here.
\end{abstract}

\maketitle

%=====================================================================================================================
\section{Introduction}
\label{sec:intro}

The particle discovered at the Large Hadron Collider in 2012~\cite{ATLAS2012,CMS2012} has, within uncertainties of a few
per cent, the couplings of the elementary Higgs boson of the Standard Model~\cite{ATLAS2022,CMS2022}. Whether the Higgs
field is elementary or a composite of more fundamental fields is nevertheless one of the oldest open questions about it.
Composite proposals range from technicolor~\cite{Weinberg1976,Susskind1979} and the condensation of the top
quark~\cite{Miransky1989,BHL1990} to a Higgs boson that is a pseudo-Goldstone boson of a strongly coupled
sector~\cite{KaplanGeorgi1984,Agashe2005}; the review of Hill and Simmons~\cite{HillSimmons2003} describes the dynamics.
Their common prototype is the model of Nambu and Jona-Lasinio~\cite{NJL1961}: a four-fermion interaction above a
critical strength condenses the fermion field, the fermions acquire a gap, and a composite scalar appears whose mass is
tied to that gap.

In earlier work~\cite{Nash2026} the author showed that the unit bi-octonions form a Moufang loop, the Lorentz super-loop,
that contains $SL(2,\mathbb C)$ together with eight spinor charges, and that a spin-zero superfluid of these charges forms
from pairs of two members of an isodoublet $\Psi_\LL$ of Dirac spinors. Such pairs are isospin triplets, so that
superfluid is an isovector; it cannot have the quantum numbers of the Higgs doublet, whatever the interaction
(Sec.~\ref{sec:model}). The present paper changes the hypotheses so that it can. Two singlet Dirac fields are added,
the hypercharges are chosen so that a pair of a member of $\Psi_\LL$ and a singlet is a doublet with $|Y|=\frac12$, and
the attraction acts in those pairs. We ask, in five mean-field methods, whether the resulting superfluid forms, which
symmetries it breaks, which gauge bosons it makes massive and with what ratio $\rho$ of their masses, how heavy its
amplitude mode is compared with the Higgs boson, and how that mode moves in time.

The answers are as follows. In every method the ground state is the custodial superfluid: it has the symmetry-breaking
pattern of the Standard Model, a massless photon and $\rho=1$, so that with its scale set to $v$ the $W$ and $Z$ bosons
have their measured masses at tree level. The two pair fields carry the scalars of two Higgs doublets. The amplitude
mode is a composite scalar, but it lies near the energy that breaks a pair, and the scale of the breaking is small
compared with that energy: the ratio $M_H/v$ is about sixteen times its measured value. Within the mean field the
superfluid is therefore the Higgs field of the Standard Model in its quantum numbers and its gauge-boson masses, but not
in the mass of its boson.

The gauge group contains $SL(2,\mathbb C)$ as a factor with a Lorentz connection, which is a pure gauge on the flat
space-time of this paper. E.~A.~Lord~\cite{Lord1971} showed that the local change of representation of the
two-component spinor algebra, made a Yang-Mills gauge group, leads with the necessary auxiliary fields to a
self-interacting spin-two field identical with Einstein's gravitation. Section~\ref{sec:gravity} outlines how that
construction will make the gravitational sector of the present model; it is a program, and no part of it is computed
here. Every number of this paper is produced by programs in Rust; independent programs in Python check the key results
with their own engine, and scripts in the Wolfram Language check the group theory and the closed forms in exact
arithmetic. The Supplemental Material~\cite{SM} gives the derivations, the checks and the data.

%=====================================================================================================================
\section{The model}
\label{sec:model}

\paragraph*{Space-time and fields.}
Space-time is the manifold M31 of the Hermitian bi-quaternions, which is Minkowski space with the metric
$\eta^\LL={\rm diag}(1,-1,-1,-1)$~\cite{Nash2026}; we use $\hbar=c=1$. There are three Dirac fields of $SL(2,\mathbb C)$,
all singlets of colour and all vector-like (both chiralities carry the same charges): the isodoublet $\Psi_\LL$, which
carries the spinor charges of the super-loop, and two singlets $\chi_\SH$ and $\chi'_\SH$. Their charges are those of
Table~\ref{tab:fields}. Every field has the mass $0.2\,\Lambda_\LL$, where $\Lambda_\LL$ is the cutoff of the momenta of
the uniform states and the unit of energy throughout.

\begin{table}[t]
\caption{The fields and their charges: the hypercharge $Y$, the electric charge $Q=T^3+Y$, the particle number
$F^{\rm tot}_\SH$ (1 for every particle of every field) and the relative number $R_\SH$ (particles of $\Psi_\LL$ minus
particles of the singlets).}
\label{tab:fields}
\begin{ruledtabular}
\begin{tabular}{lccccc}
field & $SU(2)_L$ & $Y$ & $Q$ & $F^{\rm tot}_\SH$ & $R_\SH$\\
\hline
$\Psi_\LL$ (up, down) & $\mathbf 2$ & $\frac14$ & $\frac34$, $-\frac14$ & 1 & $+1$\\
$\chi_\SH$ & $\mathbf 1$ & $\frac14$ & $\frac14$ & 1 & $-1$\\
$\chi'_\SH$ & $\mathbf 1$ & $-\frac34$ & $-\frac34$ & 1 & $-1$\\
\end{tabular}
\end{ruledtabular}
\end{table}

\paragraph*{Gauge group and covariant derivative.}
The group of the local transformations is $SL(2,\mathbb C)\times SU(2)_L\times U(1)_Y$. The isospin of $\Psi_\LL$ is
$SU(2)_L$, and on $\Psi_\LL$ and $\chi_\SH$ the hypercharge is a quarter of the fermion number of the earlier hypotheses.
The covariant derivative of each field is
\begin{equation}
D_\mu=\partial_\mu+\Gamma^\LL_\mu+ig\,\tfrac{\tau^a}{2}W^a_\mu+ig'\,Y B_\mu ,
\label{eq:cov}
\end{equation}
where $\Gamma^\LL_\mu$ is the Lorentz (spin) connection, $\tau^a/2$ acts on the isospin of $\Psi_\LL$ and is zero on the
singlets, and the couplings are those of the Standard Model: $g={{G}}$ and $g'={{GP}}$ from the measured
$G_F$, $M_W$ and $M_Z$~\cite{PDG2026}. On flat M31, $\Gamma^\LL_\mu$ is a pure gauge and is set to zero. Because every
field is vector-like, the gauge group has no anomaly.

\paragraph*{Why pairs of one doublet cannot be the Higgs field.}
The spin-zero, Lorentz-invariant bilinears of two Dirac spinors, $\psi^TC\psi$ and $\psi^TC\gamma^5\psi$, are built with
antisymmetric matrices. Fermi statistics then makes the isospin factor of a pair of two members of one doublet
symmetric, and the symmetric part of $\mathbf 2\otimes\mathbf 2$ is the triplet. A condensate of such pairs is an
isovector with the charge of two particles; no choice of the couplings changes this, because the quantum numbers of a
condensate are those of the fields that pair.

\paragraph*{Interaction.}
The attraction acts between a member of $\Psi_\LL$ and a singlet, in the spin-zero state of positive parity:
\begin{equation}
\mathcal V^{\rm pair}=-G_\SH\sum_{s={\rm u},{\rm d}}\big[(\Pi^s_\SH)^\dagger\Pi^s_\SH+(\Pi'^s_\SH)^\dagger\Pi'^s_\SH\big],
\label{eq:Vpair}
\end{equation}
with $\Pi^s_\SH=\frac12\Psi^T_\LL(C\gamma^5\otimes\mathbf e^s)\chi_\SH$ and $\Pi'^s_\SH$ the same with $\chi'_\SH$, where
$\mathbf e^{\rm u}=(1,0)^T$ and $\mathbf e^{\rm d}=(0,1)^T$. The sum is the squared length of two $SU(2)_L$ doublets and is
invariant under the gauge group. $\Psi_\LL$ keeps the scalar and the vector contact terms of the earlier hypotheses,
$\frac12G^{\rm S}_\SF{:}(\bar\Psi_\LL\Psi_\LL)^2{:}+\frac12G^{\rm V}_\SF{:}(\bar\Psi_\LL\gamma_\mu\Psi_\LL)^2{:}$; the singlets
interact only through (\ref{eq:Vpair}) and $B_\mu$. The couplings are $G^{\rm S}_\SF=-8$, $G^{\rm V}_\SF=0.8$ and
$G_\SH=100.8$ in units of $\Lambda_\LL^{-2}$; in the comparison version with one singlet the terms of $\chi'_\SH$ are
absent.

\paragraph*{Custodial symmetry and the composite field.}
Without $B_\mu$ the Lagrangian has the global symmetry $SU(2)_L\times SU(2)_R$, where $SU(2)_R$ rotates
$(\chi_\SH,\chi'_\SH)$ as a doublet (the two singlets have equal masses and one coupling), together with the relative
number $R_\SH$ and the particle number $F^{\rm tot}_\SH$. The
hypercharge is $Y=T^{3}_R+\frac14R_\SH$, so $g'$ breaks $SU(2)_R$ to its third component, as in the Standard Model, where
$SU(2)_R$ is the custodial symmetry of the Higgs field~\cite{SSVZ1980}. The pair fields form the matrix
\begin{equation}
\mathcal M_\SH=\begin{pmatrix}\Phi^{\rm u}&\Phi'^{\rm u}\\\Phi^{\rm d}&\Phi'^{\rm d}\end{pmatrix},\qquad
\Phi^s\propto\langle\Pi^s_\SH\rangle,\quad\Phi'^s\propto\langle\Pi'^s_\SH\rangle ,
\label{eq:M}
\end{equation}
which transforms as $\mathcal M_\SH\to L\,\mathcal M_\SH R^T$, a bidoublet. Its first column is a doublet with $Y=\frac12$,
the quantum numbers of the Higgs doublet; its second column has $Y=-\frac12$, and its conjugate $i\tau^2\Phi'^*$ is a
second doublet with $Y=\frac12$. The superfluid is therefore a composite with the quantum numbers of two Higgs
doublets~\cite{Branco2012}, of spin zero, positive parity, particle number two and relative number zero. A vacuum
proportional to $i\tau^2$, with the neutral components $\Phi^{\rm d}$ and $\Phi'^{\rm u}$ of equal length, keeps the
diagonal group $SU(2)_V$ and with it the electric charge $Q=T^3+T^3_R+\frac14R_\SH$; it breaks three generators of
$SU(2)_L\times U(1)_Y$ and the particle number. Its scale $v_\SH$ is defined by $M_W=\frac12gv_\SH$, as $v$ is in the
Standard Model. The model has one scale, $\Lambda_\LL$; the hypothesis that the superfluid is the Higgs field sets
$v_\SH=v=246.22$~GeV, and every other mass in GeV is then a prediction. These statements of group theory are checked in
exact arithmetic~\cite{SM}.

%=====================================================================================================================
\section{Five mean-field methods}
\label{sec:methods}

A uniform quasi-free state of the sixteen components $(\Psi_\LL,\chi_\SH,\chi'_\SH)$ is the ground state of the
Bogoliubov-de Gennes matrix
\begin{equation}
\mathsf H(\vec q)=\begin{pmatrix}\mathsf h(\vec q)+\mathsf N&\mathsf P\\\mathsf P^\dagger&-\big(\mathsf h(-\vec q)+\mathsf N\big)^T\end{pmatrix},
\label{eq:BdG}
\end{equation}
with the free Dirac matrix $\mathsf h(\vec q)$ of the three fields, a Hermitian normal self-energy $\mathsf N$ and an
antisymmetric pair self-energy $\mathsf P$. The energy of the contact terms is evaluated with Wick's theorem, and its
direct, exchange and pair contractions are kept with weights that define the methods:
\begin{itemize}
\item HS $(s_{\rm S},s_{\rm V})$: Hubbard-Stratonovich~\cite{Hubbard1959,Stratonovich1957} at the saddle point, with the
shares $s_{\rm S}$, $s_{\rm V}$ of the scalar and the vector term decoupled in density channels and the rest in pair
channels; the pairing term is decoupled in its own pair channels, and no exchange contraction is kept. We compute
$(s_{\rm S},s_{\rm V})=(1,1)$, $(0,1)$ and $(0,0)$. In HS (0,1) and HS (0,0), the pairing-only methods, only the pairing
acts on the states of this paper.
\item KS: the method of Kohn and Sham~\cite{HohenbergKohn1964,KohnSham1965}, with all three contractions, the exchange
evaluated with the projection of the local normal matrix onto the number and scalar densities of $\Psi_\LL$ and of the
pair of singlets.
\item CJT: the two-particle-irreducible effective action of Cornwall, Jackiw and Tomboulis~\cite{CJT1974}, truncated at
the contractions of the contact terms, all kept with the weight 1.
\end{itemize}
On the states computed here KS and CJT coincide; this is a theorem~\cite{SM}, not a coincidence of the numbers. The
self-energies are the derivatives of the interaction energy with respect to the local matrices; a stationary state is a
state whose self-energies are those that its own local matrices produce. The states are computed on 24 Gauss-Legendre
momenta on $[0,\Lambda_\LL]$ times the six directions of the axes, self-consistent to $10^{-13}$, by a Rust program; a
Python program with its own engine recomputes the ground states and the key results, scripts in the Wolfram Language
check the closed forms and the group theory in exact arithmetic, and a notebook repeats their checks~\cite{SM}. The time-dependent calculations use the integrator CVODE of
SUNDIALS~\cite{SUNDIALS2005,SUNDIALS2022}.

%=====================================================================================================================
\section{The ground state}
\label{sec:ground}

The normal vacuum is unstable only towards the pairs of $\Psi_\LL$ with a singlet. Its static response has the largest
eigenvalue $G_\SH/G^{\rm crit}_\SH={{LAMH}}$ in every method, eightfold (the four complex pair fields), with the critical
coupling
\begin{equation}
G^{\rm crit}_\SH=\frac{4\pi^2}{I_1(m)}={{GCRIT}}
\label{eq:Gcrit}
\end{equation}
for the mass $m=0.2\,\Lambda_\LL$, where $I_1(M)=\int_0^{\Lambda_\LL}q^2(q^2+M^2)^{-1/2}\,dq$; every other channel has
eigenvalues below 1. In every method the ground state is the
custodial superfluid: both pair fields condense with equal lengths, the down component of $\Psi_\LL$ pairs with
$\chi_\SH$ and the up component with $\chi'_\SH$, and the pair self-energy is $\mathsf P=d^0_\SH\mathsf P^{\rm vac}$ along
the direction of the vacuum $\mathsf P^{\rm vac}$ that corresponds to $\mathcal M_\SH\propto i\tau^2$. There is no pair of
two members of one field and no mixing of the fields (Table~\ref{tab:ground}). The state keeps four of the eight
generators of the global group, $SU(2)_V$ and $R_\SH$, and with them the electric charge. Its static response has exactly
four eigenvalues equal to 1, the Goldstone modes of the three broken generators of $SU(2)_L\times U(1)_Y$ and of the
particle number, and all others below 1. Every other stationary state is unstable towards it: the normal vacuum, the
condensate of one doublet (the ground state of the version with one singlet, which is stationary but has an eigenvalue
{{ONE_DOUBLET_RANGE}} in the direction of the second doublet), and, in KS and CJT, a scalar condensate without pairs in
which the Dirac masses of all four species are negative. Every constituent is paired: the lightest quasi-particle has the
energy $E^{\rm gap}$ of Table~\ref{tab:ground} (second version) and the electric charge $-\frac14$ or $\frac34$.

\begin{table}[t]
\caption{The ground states, second version (two singlets) and, for comparison, first version (one singlet): the pair
self-energy $d^0_\SH$, the energy per volume $\mathcal E$ measured from the normal vacuum, and the gap $E^{\rm gap}$, the
lowest energy of a quasi-particle; units of $\Lambda_\LL$.}
\label{tab:ground}
\begin{ruledtabular}
{{TABLE_GROUND}}
\end{ruledtabular}
\end{table}

%=====================================================================================================================
\section{The effective potential and the gauge bosons}
\label{sec:gauge}

\paragraph*{Effective potential.}
Along the direction of the vacuum, the energy of the lowest uniform state with a given pair amplitude
$\varphi_\SH=|\langle\Pi^{\rm d}_\SH\rangle|$ is the effective potential $V^{\rm eff}_\SH(\varphi_\SH)$. A Lagrange
multiplier turns the constrained state into the ground state of (\ref{eq:BdG}) with $\mathsf P=d\,\mathsf P^{\rm vac}$ for
some $d$ and a self-consistent $\mathsf N$. Because $\mathsf N$ need not be unique at fixed $d$, the states are followed by
continuation from the ground state and from the normal vacuum, and the lower branch at the same $\varphi_\SH$ is
$V^{\rm eff}_\SH$. In every method it is a Mexican hat with its minimum at the ground state (Fig.~\ref{fig:potential}). Its
top is the normal vacuum, except in KS and CJT of the second version, where it is the scalar condensate, at
$V^{\rm eff}_\SH(0)={{TOP_RATIO}}\,|V^{\rm eff}_\SH(\varphi^0_\SH)|$; there the states continued from the normal vacuum form
a second branch, above the potential, which the continuation follows up to $d={{SB_LAST}}\,d^0_\SH$.

\begin{figure}[t]
\includegraphics[width=\columnwidth]{figures/fig_prd_sh_potential.pdf}
\caption{The effective potential along the direction of the vacuum, divided by the depth of its minimum, against the
pair amplitude, for the five methods (second version). KS and CJT coincide. The dashed line with circles is the second
branch of KS and CJT, the states continued from the normal vacuum.}
\label{fig:potential}
\end{figure}

\paragraph*{Susceptibilities and masses.}
A uniform, static gauge field couples to the charge densities; its time component acts as a chemical potential
$\mu^a=g_aA^a_0$ of the charge $Q_a$. The response of the superfluid defines the susceptibilities
$\chi^{ab}=\partial\langle Q_a\rangle/\partial\mu^b$, and the effective Lagrangian
$Z_\SH|D_\mu\Phi|^2$ of the pair fields gives the gauge fields the mass matrix
\begin{equation}
(\mathsf M^{\rm g})^2_{ab}=g_ag_b\,\chi^{ab}
\label{eq:Mg}
\end{equation}
on $(W^1_\mu,W^2_\mu,W^3_\mu,B_\mu)$, with $g_a=g$ for $W^a$ and $g'$ for $B$. The normal vacuum has a gap and does not
respond, so the gauge bosons get no mass without the condensate. In the superfluid the susceptibilities are symmetric
(the theorem of Hellmann and Feynman), the row of the electric charge vanishes in every method, so the photon is
massless, and four generators are broken in the second version (three in the first). In the second version the
custodial structure
$\chi(T^a)=\chi(T^a_R)=-\chi(T^a,T^a_R)$ holds, so that
\begin{equation}
\rho_\SH=\frac{M_W^2}{M_Z^2\cos^2\theta_W}=\frac{\chi^{11}}{\chi^{33}}=1
\label{eq:rho}
\end{equation}
to {{RHO2_DEV}} in every method. With one singlet $\rho_\SH$ is {{RHO1_RANGE}}, which the measured
$M_W^2/(M_Z^2\cos^2\hat\theta_W)={{RHO_MEAS}}$~\cite{PDG2026} excludes; the second singlet is what makes the custodial
symmetry exact. In the Standard Model the departure of the measured value from 1 is a radiative correction, mainly that
of the top quark~\cite{Veltman1977}.

\paragraph*{The scale in GeV.}
From $M_W=\frac12gv_\SH$, $v_\SH=2(\chi^{11})^{1/2}$ is {{VSH2_RANGE}} in units of the cutoff (second version,
Table~\ref{tab:masses}). With $v_\SH=246.22$~GeV the cutoff is {{LAM2_GEV}}~GeV, and $M_W$ and $M_Z$ take their measured
values. The lightest fermionic excitations, which carry the fractional charges $\frac34$ and $-\frac14$, then have masses
of {{GAP_GEV_RANGE}}~GeV.

\begin{table*}[t]
\caption{The scale and the masses, second version: $v_\SH$ (units of $\Lambda_\LL$), the cutoff, $\rho_\SH$, and the
amplitude mode from the effective Lagrangian ($M^{\rm EFT}_\SH$, all methods) and from the linear response
($\omega^{\rm H}_\SH$, the pairing-only methods), relative to the scale and in GeV. The Standard Model has
$m_H/v={{MH_SM_V}}$.}
\label{tab:masses}
\begin{ruledtabular}
{{TABLE_MASSES}}
\end{ruledtabular}
\end{table*}

%=====================================================================================================================
\section{The amplitude mode}
\label{sec:amplitude}

\paragraph*{Effective Lagrangian.}
At long wavelengths the pair fields obey $Z_\SH(|\partial\Phi^{\rm d}|^2+|\partial\Phi'^{\rm u}|^2)-V^{\rm eff}_\SH$. The
kinetic coefficient follows from the phase stiffness, $\chi^{33}=Z_\SH(\varphi^0_\SH)^2$, and a common oscillation of
the two lengths has the frequency
\begin{equation}
(M^{\rm EFT}_\SH)^2=\frac{V^{\rm eff\,\prime\prime}_\SH(\varphi^0_\SH)\,(\varphi^0_\SH)^2}{4\chi^{33}} .
\label{eq:MEFT}
\end{equation}
It gives $M^{\rm EFT}_\SH/v_\SH={{MEFT2_RANGE}}$ in the five methods (Table~\ref{tab:masses}). The formula is the leading
order of an expansion in the frequency, and here the mode lies close to the energy $2E^{\rm gap}$ that breaks a pair, so
the full linear response is needed.

\paragraph*{Linear response.}
The pole of the linear response of the time-dependent mean field at zero momentum is computed in the sector of relative
number zero, where the fluctuations of the self-energies close under the symmetries of the state. At $\omega=0$ exactly
four singular values of $1-\bm\chi(\omega)\mathsf F$ vanish ($\bm\chi$ the response of the local matrices, $\mathsf F$
the kernel of the contact terms): the Goldstone modes. In HS (0,1) and HS (0,0), where only the pairing acts, the
response of the pair amplitudes takes the form of the theory of Bardeen, Cooper and Schrieffer~\cite{BCS1957}. With
$E=(q^2+M_{\rm eff}^2)^{1/2}$, $M_{\rm eff}^2=m^2+(d^0_\SH)^2/4$, the eigenvalue along the amplitude of the vacuum
direction is
\begin{equation}
\lambda^{\rm amp}(\omega)=\frac{G_\SH}{4\pi^2}\int_0^{\Lambda_\LL}\frac{q^2(q^2+m^2)\,dq}{E\,(E^2-\omega^2/4)} ,
\label{eq:lamamp}
\end{equation}
and with the gap equation $G_\SH I_1(M_{\rm eff})/(4\pi^2)=1$ the identity
\begin{equation}
\lambda^{\rm amp}(\omega)-1=\frac{\omega^2-(d^0_\SH)^2}{4}\,\frac{G_\SH}{4\pi^2}\int_0^{\Lambda_\LL}\frac{q^2\,dq}{E\,(E^2-\omega^2/4)}
\label{eq:identity}
\end{equation}
follows at every momentum. Below the threshold $2M_{\rm eff}$ every term of the integral is positive, so the amplitude
mode lies at
\begin{equation}
\omega^{\rm H}_\SH=d^0_\SH
\label{eq:omegaH}
\end{equation}
exactly, for every coupling, mass and cutoff. It is bound by $2M_{\rm eff}-d^0_\SH={{BIND2}}\,\Lambda_\LL$. This is the
analogue of the amplitude mode of a superconductor at twice the gap parameter~\cite{LittlewoodVarma1981,PekkerVarma2015};
here the mass of the fields raises the threshold above the mode. The Rust program, which solves the full response of
the sector and does not use (\ref{eq:lamamp}), finds (\ref{eq:omegaH}) to {{OMEGA_D0_DEV}}; the Python program confirms
(\ref{eq:lamamp}) numerically, and the Wolfram script derives (\ref{eq:identity}) symbolically. With the scale of the
second version, $M_H/v_\SH={{MH_RPA_V}}$, that is $M_H={{MH_RPA_GEV}}$~GeV, against $m_H/v={{MH_SM_V}}$: the amplitude
mode is {{FACTOR_RPA}} times too heavy relative to the scale.

\paragraph*{No coupling repairs it.}
In HS (0,0) the ground state has a closed form, so the ratio can be followed over the coupling. For
$1.02\le G_\SH/G^{\rm crit}_\SH\le3$ and the masses $0.1$ and $0.2$ of the fields it lies between {{SCAN_MIN2}} and
{{SCAN_MAX2}} in the second version (Fig.~\ref{fig:scan}), where $\rho_\SH=1$ at every point. The ratio is smallest near
the critical coupling,
where $d^0_\SH$ and $v_\SH$ both vanish, and it grows with the coupling, because $d^0_\SH$ grows while $v_\SH$, limited by
the cutoff, stays below {{SCAN_VMAX}}. Even its smallest value is {{SCAN_MIN_FACTOR}} times the measured $m_H/v$.

\begin{figure}[t]
\includegraphics[width=\columnwidth]{figures/fig_prd_sh_scan.pdf}
\caption{The ratio of the mass of the amplitude mode to the scale, $\omega^{\rm H}_\SH/v_\SH$, against the coupling, in
HS (0,0) for the masses 0.1 and 0.2 of the fields, in both versions, with the measured $m_H/v$.}
\label{fig:scan}
\end{figure}

\paragraph*{The partners of the amplitude mode.}
On all eight real components of the bidoublet the response of HS (0,0) has four modes at $\omega=0$ and four at
$\omega=d^0_\SH$. The generators of $SU(2)_V$ commute with the response and split each set into a singlet and a triplet.
At $\omega=0$ the triplet becomes the longitudinal parts of $W^\pm$ and $Z$, and the singlet is the Goldstone boson of the
particle number, a massless neutral scalar that the model has to remove, by gauging $F^{\rm tot}_\SH$ or by breaking it
explicitly; this is left open. At $\omega=d^0_\SH$ the singlet is the candidate for the Higgs boson, and the triplet holds a
neutral scalar and a charged scalar with its antiparticle, degenerate with it in the pairing-only methods: the scalar
content of two Higgs doublets.

%=====================================================================================================================
\section{Dynamics in time}
\label{sec:dynamics}

The time-dependent mean field of each method is integrated with CVODE (relative tolerance $10^{-12}$). The energy, the
charges of the symmetries and the orthonormality of the quasi-particle vectors are conserved to the accuracy of the
integration, and the ground state stays at rest. A small displacement of the pair fields along the vacuum direction
oscillates, in HS (0,1) and HS (0,0), at $d^0_\SH$ to {{DY_DEV_AMP}}, without damping, in both versions; a
displacement along the custodial triplet oscillates at the same frequency. In HS (1,1) the mode lies at
{{DY_RATIO_HS11}}\,$d^0_\SH$, below the energy that breaks a pair, and is undamped; the scalar and the vector terms raise
the triplet by {{DY_SPLIT_HS11}} per cent. In KS and CJT the spectrum peaks at {{DY_AMP_KS}}\,$\Lambda_\LL$, above the
threshold {{DY_GAP2_KS}}\,$\Lambda_\LL$, with a width {{DY_WR_KS}} times that of an undamped oscillation of the same
length: there the amplitude mode decays into pairs of quasi-particles (Fig.~\ref{fig:modes}). With a chemical potential
$\mu$ of the particle number the phase of the pair fields turns at the rate $2\mu$, the Josephson relation, to
{{DY_DEV_JOS}}. A displacement by 30 per cent in HS (0,0) oscillates at {{DY_LARGE_RATIO}}\,$d^0_\SH$: the oscillation
softens with its amplitude, as on a Mexican hat.

\begin{figure}[t]
\includegraphics[width=\columnwidth]{figures/fig_prd_sh_modes.pdf}
\caption{The spectra of the length $|d(t)|$ of the pair self-energy after a small displacement along the vacuum
direction, second version, each divided by its largest value; $d^0_\SH$ of the pairing-only methods is marked. KS and
CJT coincide.}
\label{fig:modes}
\end{figure}

%=====================================================================================================================
\section{Can the superfluid be the Higgs field of the Standard Model?}
\label{sec:higgs}

\paragraph*{What works.}
The superfluid of the second version has the quantum numbers of the Higgs field, breaks $SU(2)_L\times U(1)_Y$ to the
$U(1)$ of the electric charge, keeps the photon massless and gives $\rho_\SH=1$. With the scale set to $v$, the $W$ and
$Z$ bosons have the masses of the Standard Model at tree level. These are consequences of the quantum numbers and of the
custodial vacuum, and they hold in all five methods.

\paragraph*{What fails.}
The amplitude mode is {{FACTOR_RANGE}} times too heavy relative to the scale ({{FACTOR_RPA}} times from the linear
response, {{FACTOR_EFT}} times from the effective Lagrangian). In a superfluid of fermion pairs bound by a
contact interaction the amplitude mode lies at $d^0_\SH$, just below the energy $2M_{\rm eff}$ that breaks a pair, while
$v_\SH^2$, the phase stiffness, grows only like $M_{\rm eff}^2$ times the logarithm of $\Lambda_\LL^2/M_{\rm eff}^2$ (in the
leading approximation, divided by $4\pi^2$ and multiplied by the number of paired species). Here the cutoff is only about
twice $M_{\rm eff}$, so $v_\SH$ is small compared with $M_{\rm eff}$ and the ratio is large. The same relation underlies the
models in which the Higgs doublet is a bound state of the top quark and its antiparticle, where a light Higgs boson
needs a very high cutoff~\cite{BHL1990}.

\paragraph*{What the model adds.}
The second doublet brings a custodial triplet of massive scalars, degenerate with the candidate in the pairing-only
methods and split from it by {{DY_SPLIT_HS11}} per cent in HS (1,1); the particle number brings a massless Goldstone boson; and the
constituents are fermions with the fractional electric charges $\frac34$ and $-\frac14$ and masses of
{{GAP_GEV_RANGE}}~GeV. Each of these must be confronted with experiment before the model can be compared with the
observed Higgs sector.

\paragraph*{What remains open.}
The mean field leaves out the fluctuations of the pair fields, which shift both the stiffness and the amplitude mode; a
computation beyond it is needed before the ratio $M_H/v$ can be called a prediction of the model. The thermal transition
of the revised model has not been computed; the width of its critical region, which the earlier hypotheses gave as wider
than the critical temperature itself, is therefore not known here.

%=====================================================================================================================
\section{Is the Higgs a composite particle?}
\label{sec:composite}

In this model the Higgs field has no elementary scalar: its order parameter is the expectation value of the pair
operators, a composite of a member of $\Psi_\LL$ and a singlet. The amplitude mode is a bound state of two
quasi-particles below the threshold that breaks a pair in HS (0,1), HS (0,0) and HS (1,1), and a resonance above it,
decaying into pairs, in KS and CJT. The Higgs particle of the model is therefore composite by construction, and whether
it is a sharp state depends on the method; in none of them is it light enough to be the observed Higgs boson.

%=====================================================================================================================
\section{Toward the gravitational sector}
\label{sec:gravity}

The factor $SL(2,\mathbb C)$ of the gauge group acts on the spinor indices of all three fields with the Lorentz connection
$\Gamma^\LL_\mu$ of (\ref{eq:cov}), which on flat M31 is a pure gauge. Lord~\cite{Lord1971} considered the group of
changes of representation of the two-component spinor algebra, under which the free field equations of particles of any
spin are invariant, and made it a gauge group in the sense of Yang and Mills. Gauge invariance and Lorentz covariance
require auxiliary fields; Lord interpreted them as the potentials of a spin-two field, which is then self-interacting,
and showed that the Bargmann-Wigner equations of the linear spin-two field, modified by this self-interaction, give a
nonlinear theory identical with Einstein's. This construction is the natural route for the present model because its
spinor charges come with $SL(2,\mathbb C)$ from the super-loop~\cite{Nash2026}: the program is (i) to replace the
pure-gauge $\Gamma^\LL_\mu$ by Lord's gauge field and its auxiliary fields, (ii) to couple the spinor fields and their
pairs to it, and (iii) to obtain the gravitational field equations with the energy-momentum tensor of the spinor fields
as their source. Gauging the local Lorentz group is also the route of Utiyama and Kibble~\cite{Utiyama1956,Kibble1961,Hehl1976}.
The superfluid is a Lorentz scalar and leaves $SL(2,\mathbb C)$ unbroken, so the condensation does not give the
gravitational field a mass. The quantization of the gravitational sector, and with it the ``quantum gravity'' of the
title, is beyond this paper.

%=====================================================================================================================
\section{Conclusions}
\label{sec:conclusions}

Two singlet Dirac fields with the hypercharges $\frac14$ and $-\frac34$, added to the isodoublet of spinor charges of the
Lorentz super-loop, turn its superfluid into a composite with the quantum numbers of two Higgs doublets. In five
mean-field methods its ground state is the custodial vacuum, which breaks the electroweak group to electromagnetism,
gives $\rho=1$, and, with its scale set to $v$, gives the $W$ and $Z$ bosons their masses. Its amplitude mode, a composite
scalar, lies exactly at the pair self-energy in the methods in which only the pairing acts, and is about sixteen times
too heavy relative to the scale. The model therefore reproduces the Higgs field of the Standard Model in its quantum
numbers and its gauge-boson masses, but, in the mean field, not in the mass of its boson. The gravitational sector will
be added through Lord's gauge theory of the local change of spinor representation, for which the model already has the
place.

%=====================================================================================================================
\section*{Data availability}

The data that support the findings of this article are openly available in the repository of PROJECT
spacetimes~\cite{Repo}, folder \texttt{Revision\_SM\_Higgs}: the data files (\texttt{data/}), the Rust, Python and Wolfram
Language programs that produce and check them (\texttt{rust/}, \texttt{python/}, \texttt{wolfram/}), the reports of their
checks, the notebook, and the scripts that build this article from the data (\texttt{paper\_prd/}). The Supplemental
Material~\cite{SM} lists the files with their SHA-256 checksums.

\begin{acknowledgments}
The computations used the SUNDIALS integrators in a Rust implementation, numpy and LAPACK.
\end{acknowledgments}

\bibliography{references}

\end{document}
